\section{Semilattices and enriched categories}
\label{app:background}

This appendix recalls the semilattice and enriched-category background used in the main text.

We use the following standard monoidal structure on semilattices throughout.

\begin{proposition}
\label{prop:slat_closed_smc}
The category $\catname{SLat}$ of semilattices and semilattice homomorphisms is symmetric monoidal closed, with tensor product determined by maps that preserve joins separately in each variable.
\end{proposition}

\subsection{Enriched categories and semilattices}
Enriched categories were introduced by~Eilenberg and Kelly~\cite{EilenbergKellyClosedCategories} as a generalisation of ordinary categories in which hom-sets are replaced by objects of a monoidal category. 
This framework allows one to encode additional structure directly into the notion of morphism, making it possible to treat contexts such as topology, algebra, and order theory in a unified way. 
A key motivating example, due to~Lawvere~\cite{lawvere_metric_1973}, shows that metric spaces can be viewed as categories enriched over $[0,\infty]$, illustrating how familiar mathematical structures arise naturally in this setting.
This perspective not only provided new insights into category theory by drawing on intuition from metric spaces, but also allowed metric spaces themselves to be studied using categorical methods. 
As a result, enriched category theory provides a flexible and powerful language for organizing and comparing diverse mathematical phenomena.
In this work, our goals are more utilitarian, as we use the framework of enriched categories to describe a semantics of e-graphs and to define a type theory for them.
We will proceed with the recap of enriched category theory.

\begin{definition}[Definition 6.2.1 in~\cite{Borceux_1994}]
  Let $\mathcal{V} = (\mathcal{V}_{0}, \I_{\mathcal{V}}, \otimes, \lambda, \rho, \alpha)$ be a monoidal category where $\lambda$ and $\rho$ are left and right unitors and $\alpha$ is the associator.
  We define a (small) $\mathcal{V}$-enriched category $\enriched{C}$ consisting of the following data
  \begin{enumerate}
    \item a set of objects $\obj{\enriched{C}}$
    \item for every pair of objects $A,B \in \obj{\enriched{C}}$---an (hom-)object $\obj{\enriched{C}}(A,B) \in \mathcal{V}$
    \item for every triple of objects $A,B,C \in \obj{\enriched{C}}$---a composition morphism in $\mathcal{V}$
    \[
    \circ_{A,B,C} : \enriched{C}(B,C) \otimes \enriched{C}(A,B) \to \enriched{C}(A,C)
    \]
    \item for every object $A \in \obj{\enriched{C}}$---a unit morphism
    \[
      \id_{A} : \I_{\mathcal{V}} \to \enriched{C}(A,A)
    \]
  \end{enumerate}
  such that the following diagrams commute (associativity and unitality)
  % https://q.uiver.app/#q=WzAsNSxbMCwwLCIoXFxlbnJpY2hlZHtDfShDLEQpIFxcb3RpbWVzIFxcZW5yaWNoZWR7Q30oQixDKSkgXFxvdGltZXMgXFxlbnJpY2hlZHtDfShBLEIpIl0sWzAsMiwiXFxlbnJpY2hlZHtDfShDLEQpIFxcb3RpbWVzIChcXGVucmljaGVke0N9KEIsQykpIFxcb3RpbWVzIFxcZW5yaWNoZWR7Q30oQSxCKSJdLFsyLDAsIlxcZW5yaWNoZWR7Q30oQixEKSBcXG90aW1lcyBcXGVucmljaGVke0N9KEEsQikiXSxbMiwyLCJcXGVucmljaGVke0N9KEMsRClcXG90aW1lc1xcZW5yaWNoZWR7Q30oQSxDKSJdLFszLDEsIlxcZW5yaWNoZWR7Q30oQSxEKSJdLFswLDIsIlxcY2lyY197QixDLER9IFxcb3RpbWVzIFxcaWQiXSxbMCwxLCJcXGFscGhhIiwyXSxbMSwzLCJcXGlkIFxcb3RpbWVzIFxcY2lyY197QSxCLEN9IiwyXSxbMiw0LCJcXGNpcmNfe0EsQixEfSJdLFszLDQsIlxcY2lyY197QSxDLER9IiwyXV0=
\[\begin{tikzcd}
	{(\enriched{C}(C,D) \otimes \enriched{C}(B,C)) \otimes \enriched{C}(A,B)} && {\enriched{C}(B,D) \otimes \enriched{C}(A,B)} & \\
	&&& {\enriched{C}(A,D)} \\
	{\enriched{C}(C,D) \otimes (\enriched{C}(B,C)) \otimes \enriched{C}(A,B)} && {\enriched{C}(C,D)\otimes\enriched{C}(A,C)}
	\arrow["{\circ_{B,C,D} \otimes \id_{\enriched{C}(A,B)}}", from=1-1, to=1-3]
	\arrow["\alpha"', from=1-1, to=3-1]
	\arrow["{\circ_{A,B,D}}", from=1-3, to=2-4]
	\arrow["{\id_{\enriched{C}(C,D)} \otimes \circ_{A,B,C}}"', from=3-1, to=3-3]
	\arrow["{\circ_{A,C,D}}"', from=3-3, to=2-4]
\end{tikzcd}\]
% https://q.uiver.app/#q=WzAsNixbMCwxLCJcXGVucmljaGVke0N9KEEsQikiXSxbMSwwLCJcXGVucmljaGVke0N9KEEsQikgXFxvdGltZXMgXFxJX3tcXG1hdGhjYWx7Vn19Il0sWzEsMiwiXFxJX3tcXG1hdGhjYWx7Vn19IFxcb3RpbWVzIFxcZW5yaWNoZWR7Q30oQSxCKSJdLFszLDAsIlxcZW5yaWNoZWR7Q30oQSxCKSBcXG90aW1lcyBcXGVucmljaGVke0EsQX0iXSxbMywyLCJcXGVucmljaGVke0N9KEIsQikgXFxvdGltZXMgXFxlbnJpY2hlZHtDfShBLEIpIl0sWzQsMSwiXFxlbnJpY2hlZHtDfShBLEIpIl0sWzAsMSwiXFxyaG9eey0xfSIsMl0sWzAsMiwiXFxsYW1iZGFeey0xfSJdLFsxLDMsIlxcaWRfe1xcZW5yaWNoZWR7Q30oQSxCKX0gXFxvdGltZXMgXFxpZF97QX0iXSxbMiw0LCJcXGlkX3tCfSBcXG90aW1lcyBcXGlkX3tcXGVucmljaGVke0N9KEEsQil9IiwyXSxbMyw1LCJcXGNpcmNfe0EsQSxCfSIsMl0sWzQsNSwiXFxjaXJjX3tBLEIsQn0iXV0=
\[\begin{tikzcd}
	& {\enriched{C}(A,B) \otimes \I_{\mathcal{V}}} && {\enriched{C}(A,B) \otimes \enriched{C}(A,A)} & \\
	{\enriched{C}(A,B)} &&&& {\enriched{C}(A,B)} \\
	& {\I_{\mathcal{V}} \otimes \enriched{C}(A,B)} && {\enriched{C}(B,B) \otimes \enriched{C}(A,B)}
	\arrow["{\id_{\enriched{C}(A,B)} \otimes \id_{A}}", from=1-2, to=1-4]
	\arrow["{\circ_{A,A,B}}"', from=1-4, to=2-5]
	\arrow["{\rho^{-1}}"', from=2-1, to=1-2]
	\arrow["{\lambda^{-1}}", from=2-1, to=3-2]
	\arrow["{\id_{B} \otimes \id_{\enriched{C}(A,B)}}"', from=3-2, to=3-4]
	\arrow["{\circ_{A,B,B}}", from=3-4, to=2-5]
\end{tikzcd}\]
\end{definition}

We will then omit writing the subscripts, when they are clear from the context.
We will also replace the composition of $\lambda, \rho, \alpha$ by writing $\cong$.

As hom-objects are not generally sets, we cannot talk about elements of them.
Instead, we can talk about morphisms from the unit $\I_{\mathcal{V}}$ to a given hom-object $\enriched{C}(A,B)$, which are called \emph{generalised elements} of $\enriched{C}(A,B)$, like we saw above in the case of $\id_{A}$.
For example, the notion of isomorphism between two objects $A$ and $B$ in such a category $\enriched{C}$ takes the following form.

\begin{definition}[Isomorphism in $\enriched{C}$. Lemma 3.5.12~\cite{Riehl_2014}]

Let $\enriched{C}$ be a $\mathcal{V}$-enriched category.
Given two objects $A,B \in \enriched{C}$ and a morphism $f : I_{\mathcal{V}} \to \enriched{C}(A,B)$, we call $f$ an isomorphism if there is a $g : I_{\mathcal{V}} \to \enriched{C}(B,A)$ such that
% https://q.uiver.app/#q=WzAsNCxbMCwwLCJcXEkiXSxbMiwwLCJcXEkgXFxvdGltZXMgXFxJIl0sWzQsMCwiXFxlbnJpY2hlZHtDfShCLEEpIFxcb3RpbWVzIFxcZW5yaWNoZWR7Q30oQSxCKSJdLFs0LDIsIlxcZW5yaWNoZWR7Q30oQSxBKSJdLFsyLDMsIlxcY2lyYyJdLFswLDMsIlxcaWRfe0F9IiwyXSxbMCwxLCJcXGNvbmciXSxbMSwyLCJmIFxcb3RpbWVzIGciXV0=
\[\begin{tikzcd}
	\I && {\I \otimes \I} && {\enriched{C}(B,A) \otimes \enriched{C}(A,B)} \\
	\\
	&&&& {\enriched{C}(A,A)}
	\arrow["\cong", from=1-1, to=1-3]
	\arrow["{\id_{A}}"', from=1-1, to=3-5]
	\arrow["{f \otimes g}", from=1-3, to=1-5]
	\arrow["\circ", from=1-5, to=3-5]
\end{tikzcd}\]
\end{definition}

\begin{example}
    Since $\catname{SLat}$ is a set-like category, we can talk about elements of the hom-objects for an $\catname{SLat}$-category $\enriched{C}$ by pretending that a particular morphism was chosen by a map from the unit semilattice.
    Thus, we will simply write $f \in \enriched{C}(A,B)$ meaning $\I_{\catname{SLat}} \to \enriched{C}(A,B)$.
    $\circ$ being a homomorphism of semilattices means that for any $f,f' \in \enriched{C}(A,B)$ and $g,g' \in \enriched{C}(B,C)$ we have
    \[
    (f' \join f) \circ (g' \join g) = (f' \circ g') \join (f' \circ g) \join (f \circ g') \join (f \circ g)
    \]
\end{example}

\begin{definition}
  Let $\enriched{C}$ and $\enriched{D}$ be two $\mathcal{V}$-categories.
  A $\mathcal{V}$-enriched functor $F : \enriched{C} \to \enriched{D}$ is defined by the following data.
  \begin{enumerate}
    \item A mapping $F : \obj{\enriched{C}} \to \obj{\enriched{D}}$
    \item An object-indexed family of morphisms in $\mathcal{V}$ $F_{A,B} : \enriched{C}(A,B) \to \enriched{D}(FA,FB)$
  \end{enumerate}
  such that the following diagrams commute

  \[
  \begin{tikzcd}
	{\enriched{C}(A,A') \otimes \enriched{C}(A',A'')} && {\enriched{C}(A,A'')} \\
	\\
	{\enriched{D}(FA,FA') \otimes \enriched{D}(FA',FA'')} && {\enriched{D}(FA,FA'')}
	\arrow["{c_{A,A',A''}}"', from=1-1, to=1-3]
	\arrow["{F_{A,A'} \otimes F_{A',A''}}", from=1-1, to=3-1]
	\arrow["{F_{A,A''}}"', from=1-3, to=3-3]
	\arrow["{c_{FA,FA',FA''}}"', from=3-1, to=3-3]
\end{tikzcd}
\]

\[
\begin{tikzcd}
	I_{\mathcal{V}} && {\enriched{C}(A,A)} \\
	\\
	&& {\enriched{D}(FA,FA)}
	\arrow["{u_{A}}", from=1-1, to=1-3]
	\arrow["{u_{FA}}"', from=1-1, to=3-3]
	\arrow["{F_{A,A}}", from=1-3, to=3-3]
\end{tikzcd}
\]
\end{definition}

\begin{definition}
  Given two $\mathcal{V}$-functors $F,G : \enriched{C} \to \enriched{D}$ a $\mathcal{V}$-natural transformation is family of morphisms indexed by objects in $\enriched{C}$, $\alpha_{A} : I_{\mathcal{V}} \to \enriched{D}(FA,GA)$ in $\mathcal{V}$ such that the following diagram commutes
  \[
  \begin{tikzcd}
	& {\enriched{C}(A,A')\otimes I_{\mathcal{V}}} && {\enriched{D}(FA,FA') \otimes \enriched{D}(FA',GA')} \\
	{\enriched{C}(A,A')} &&&& {\enriched{D}(FA,GA')} \\
	& {I_{\mathcal{V}} \otimes \enriched{C}(A,A')} && {\enriched{D}(FA,GA) \otimes \enriched{D}(GA,GA')}
	\arrow["{G_{A,A'}\otimes \alpha_{A'}}"', from=1-2, to=1-4]
	\arrow["{c_{FA,FA',GA'}}"', from=1-4, to=2-5]
	\arrow["{r^{-1}}", from=2-1, to=1-2]
	\arrow["{l^{-1}}", from=2-1, to=3-2]
	\arrow["{\alpha_{A} \otimes F_{A,A'}}", from=3-2, to=3-4]
	\arrow["{c_{FA,GA,GA'}}", from=3-4, to=2-5]
\end{tikzcd}
\]
\end{definition}

The data above defines a 2-category $\mathcal{V}\text{-}\catname{Cat}$ of $\mathcal{V}$-enriched categories with $\mathcal{V}$-functors being 1-cells and $\mathcal{V}$-natural transformations being 2-cells.
This in particular means that we can define adjunction as in any 2-category which immediately gives the definition of many standard categorical constructions.

\begin{definition}[Page 12 in~Kelly et al.~\cite{Kelly2022BASICCO}]
    Let $\mathcal{V}$ by a symmetric monoidal category.
    Then, for any $\mathcal{V}$-categories $\mathcal{M}$ and $\mathcal{N}$ we can define the product category $\mathcal{M} \times \mathcal{N}$ defined by the following data.
    \paragraph{Objects.} $\obj{\mathcal{M} \times \mathcal{N}} = \obj{\mathcal{M}} \times \obj{\mathcal{N}}$.
    \paragraph{Hom-objects.} $\mathcal{M} \times \mathcal{N}((A,B), (C,D)) = \mathcal{M}(A,C) \otimes \mathcal{N}(B,D)$.
    The composition is given be the top arrow in
    % https://q.uiver.app/#q=WzAsMyxbMCwwLCIoXFxtYXRoY2Fse019KEIsRSkgXFxvdGltZXMgXFxtYXRoY2Fse059KEQsRikpIFxcb3RpbWVzIChcXG1hdGhjYWx7TX0oQSxCKSBcXG90aW1lcyBcXG1hdGhjYWx7Tn0oQyxEKSkiXSxbMCwyLCIoXFxtYXRoY2Fse019KEIsRSkgXFxvdGltZXMgXFxtYXRoY2Fse019KEEsQikpIFxcb3RpbWVzIChcXG1hdGhjYWx7Tn0oRCxGKSBcXG90aW1lcyBcXG1hdGhjYWx7Tn0oQyxEKSkiXSxbMiwwLCJcXG1hdGhjYWx7TX0oQSxFKSBcXG90aW1lcyBcXG1hdGhjYWx7Tn0oQyxGKSJdLFswLDEsIm0iLDJdLFsxLDIsIlxcY2lyY197XFxtYXRoY2Fse019fSBcXG90aW1lcyBcXGNpcmNfe1xcbWF0aGNhbHtOfX0iLDJdLFswLDIsIlxcY2lyY197XFxtYXRoY2Fse019IFxcdGltZXMgXFxtYXRoY2Fse059fSJdXQ==
\[\begin{tikzcd}
	{(\mathcal{M}(B,E) \otimes \mathcal{N}(D,F)) \otimes (\mathcal{M}(A,B) \otimes \mathcal{N}(C,D))} && {\mathcal{M}(A,E) \otimes \mathcal{N}(C,F)} \\
	\\
	{(\mathcal{M}(B,E) \otimes \mathcal{M}(A,B)) \otimes (\mathcal{N}(D,F) \otimes \mathcal{N}(C,D))}
	\arrow["{\circ_{\mathcal{M} \times \mathcal{N}}}", from=1-1, to=1-3]
	\arrow["m"', from=1-1, to=3-1]
	\arrow["{\circ_{\mathcal{M}} \otimes \circ_{\mathcal{N}}}"', from=3-1, to=1-3]
\end{tikzcd}
\]
where $m$ is any suitable composition of $\alpha$ and $\sym$.
The identity is given by the morphism
\[
\I \xrightarrow{\cong} \I \otimes \I \xrightarrow{\id_{\mathcal{M}} \otimes \id_{\mathcal{N}}} \mathcal{M}(A,A) \otimes \mathcal{N}(B,B)
\]
\end{definition}

\begin{remark}
    Note that symmetry of $\mathcal{V}$ is essential to define a tensor product in a $\mathcal{V}$-category.
\end{remark}

\begin{definition}
\label{def:monoidal_enriched_category}
    Let $\enriched{C}$ be a $\mathcal{V}$-category.
    $\enriched{C}$ is a monoidal $\mathcal{V}$-category if it is equipped with the following functors
    \[
    \enriched{C} \times \enriched{C} \xrightarrow{\enriched{\otimes}} \enriched{C} \qquad \enriched{I} \xrightarrow{\I} \enriched{C}
    \]
    together with appropriate associators and unitors, where $\enriched{I}$ is the unit $\mathcal{V}$-category with one object and hom-object given by $\I_{\mathcal{V}}$.
    The definition of composition in $\enriched{C} \times \enriched{C}$ together with functoriality of $\enriched{\otimes}$ results in the following diagram
    % https://q.uiver.app/#q=WzAsNCxbMCwwLCIoXFxlbnJpY2hlZHtDfShCLEUpIFxcb3RpbWVzIFxcZW5yaWNoZWR7Q30oRCxGKSkgXFxvdGltZXMgKFxcZW5yaWNoZWR7Q30oQSxCKSBcXG90aW1lcyBcXGVucmljaGVke0N9KEMsRCkpIl0sWzMsMCwiXFxlbnJpY2hlZHtDfShBLEUpIFxcb3RpbWVzIFxcZW5yaWNoZWR7Q30oQyxGKSJdLFswLDIsIlxcZW5yaWNoZWR7Q30oQiBcXGVucmljaGVke1xcb3RpbWVzfSBELCBFIFxcZW5yaWNoZWR7XFxvdGltZXN9IEYpIFxcb3RpbWVzIFxcZW5yaWNoZWR7Q30oQSBcXGVucmljaGVke1xcb3RpbWVzfSBDLCBCIFxcZW5yaWNoZWR7XFxvdGltZXN9IEQpIl0sWzMsMiwiXFxlbnJpY2hlZHtDfShBIFxcZW5yaWNoZWR7XFxvdGltZXN9IEMsIEUgXFxlbnJpY2hlZHtcXG90aW1lc30gRikiXSxbMCwxLCJcXGNpcmN7XFxlbnJpY2hlZHtDfSBcXHRpbWVzIFxcZW5yaWNoZWR7Q319Il0sWzAsMiwiXFxlbnJpY2hlZHtcXG90aW1lc30gXFxvdGltZXMgXFxlbnJpY2hlZHtcXG90aW1lc30iLDJdLFsxLDMsIlxcZW5yaWNoZWR7XFxvdGltZXN9Il0sWzIsMywiXFxjaXJjX3tcXGVucmljaGVke0N9fSIsMl1d
\[\begin{tikzcd}
	{(\enriched{C}(B,E) \otimes \enriched{C}(D,F)) \otimes (\enriched{C}(A,B) \otimes \enriched{C}(C,D))} &&& {\enriched{C}(A,E) \otimes \enriched{C}(C,F)} \\
	\\
	{\enriched{C}(B \enriched{\otimes} D, E \enriched{\otimes} F) \otimes \enriched{C}(A \enriched{\otimes} C, B \enriched{\otimes} D)} &&& {\enriched{C}(A \enriched{\otimes} C, E \enriched{\otimes} F)}
	\arrow["{\circ_{\enriched{C} \times \enriched{C}}}", from=1-1, to=1-4]
	\arrow["{\enriched{\otimes} \otimes \enriched{\otimes}}"', from=1-1, to=3-1]
	\arrow["{\enriched{\otimes}}", from=1-4, to=3-4]
	\arrow["{\circ_{\enriched{C}}}"', from=3-1, to=3-4]
\end{tikzcd}\]
meaning that $\enriched{\otimes}$ distributes over $\otimes$.
\end{definition}

\begin{example}
    For an $\catname{SLat}$-category $\enriched{C}$, the hom-objects of $\enriched{C} \times \enriched{C}$ are tensor products of semilattices.
    This, together with the tensor being a homomorphism of semilattices, means that we have the following property for the tensor product
    \[
    (f \join g) \otimes h = (f \otimes h) \join (g \otimes h),
    \]
    because in a tensor product of semilattices we have
    $(f \join g,h)=(f,h)\join(g,h)$, so the action of the tensor-product
    functor must agree on both sides.
\end{example}

\begin{example}
    We can then define a closed monoidal $\catname{SLat}$-category $\enriched{C}$ as a category equipped with a pair of adjoint $\catname{SLat}$-functors for all objects $X$
    \[
    (-) \otimes X \dashv X \multimap (-) \colon \enriched{C} \to \enriched{C}
    \]
    defined by $\catname{SLat}$-natural transformations
    \[
    \eta \colon \I \to X \multimap (- \otimes X) \qquad \varepsilon \colon (X \multimap -) \otimes X.
    \]
    With these we can define the \enquote{currying} operator $\Lambda : \enriched{C}(A \otimes B, D) \to \enriched{C}(A, B \multimap D)$ for a morphism $f \in \enriched{C}(A \otimes B, D)$ as
    \[
    \Lambda(f) \Coloneq (B \multimap f) \circ \eta_{B}.
    \]
    We have that
    \begin{align*}
    \Lambda(f \join g) &= (B \multimap (f\join g)) \circ \eta_{B}\\
                   &= ((B \multimap f) \join (B \multimap g)) \circ \eta_{B}\\
                   &= ((B \multimap f) \circ \eta_{B}) \join ((B \multimap g) \circ \eta_{B})\\
                   &= \Lambda(f) \join \Lambda(g)
    \end{align*}
    by using the properties of $B \multimap -$ being an $\catname{SLat}$-functor and $\circ$ being the composition.
\end{example}

It comes as no surprise that we can turn a given ordinary category $\mathcal{C}$ into an $\catname{SLat}$-category $\enriched{C}$ is a canonical way.
This fact is leveraged when defining the semantics of e-graphs since the latter are defined over a signature of generating morphisms.
Canonically building an $\catname{SLat}$-category from this signature induces the mentioned semantics.

\begin{definition}[Proposition 6.4.7~\cite{Borceux_1994}]
    Let $\mathcal{V}$ and $\mathcal{W}$ be monoidal categories and let $F : \mathcal{V} \to \mathcal{W}$ be a monoidal functor.
    It induces a 2-functor $\tilde{F} : \mathcal{V}\text{-}\catname{Cat} \to \mathcal{W}\text{-}\catname{Cat}$ defined as follows.
    \paragraph{Objects.} Given a $\mathcal{V}$-category $\enriched{C}$, $\tilde{F}(\enriched{C})$ has the same objects as $\enriched{C}$.
    \paragraph{Hom-objects.} For any pair of objects $A,B$, $\tilde{F}(\enriched{C})(A,B) \Coloneq F(\enriched{C}(A,B))$.
    \paragraph{Composition.} For any triple of objects $A,B,C$, the composition morphism $\circ_{A,B,C} : \tilde{F}(\enriched{C})(B,C) \otimes \tilde{F}(\enriched{C})(A,B) \to \tilde{F}(\enriched{C})(A,C)$ is given by the following composition
    % https://q.uiver.app/#q=WzAsNCxbMCwwLCJcXHRpbGRle0Z9KFxcZW5yaWNoZWR7Q30pKEIsQykgXFxvdGltZXMgXFx0aWxkZXtGfShcXGVucmljaGVke0N9KShBLEIpIl0sWzAsMiwiRihcXGVucmljaGVke0N9KEIsQykpIFxcb3RpbWVzIEYoXFxlbnJpY2hlZHtDfShBLEIpKSJdLFsyLDIsIkYoXFxlbnJpY2hlZHtDfShCLEMpIFxcb3RpbWVzIFxcZW5yaWNoZWR7Q30oQSxCKSkiXSxbNCwyLCJGKFxcZW5yaWNoZWR7Q30oQSxDKSkgPSBcXHRpbGRle0Z9KFxcZW5yaWNoZWR7Q30pKEEsQykiXSxbMCwxLCI9IiwyXSxbMSwyLCJcXGNvbmciXSxbMiwzLCJGKFxcY2lyYykiXV0=
\[\begin{tikzcd}
	{\tilde{F}(\enriched{C})(B,C) \otimes \tilde{F}(\enriched{C})(A,B)} &&&& \\
	\\
	{F(\enriched{C}(B,C)) \otimes F(\enriched{C}(A,B))} && {F(\enriched{C}(B,C) \otimes \enriched{C}(A,B))} && {F(\enriched{C}(A,C)) = \tilde{F}(\enriched{C})(A,C)}
	\arrow["{=}"', from=1-1, to=3-1]
	\arrow["\cong", from=3-1, to=3-3]
	\arrow["{F(\circ)}", from=3-3, to=3-5]
\end{tikzcd}\]
\paragraph{Identities.} For any object $A$, the identity morphism $\id_{A} : I_{\mathcal{W}} \to \tilde{F}(\enriched{C})(A,A)$ is given by the following composition
\[
\I_{\mathcal{W}} \xrightarrow{\cong} F(\I_{\mathcal{V}}) \xrightarrow{F(\id_{A})} F(\enriched{C}(A,A)) = \tilde{F}(\enriched{C})(A,A)
\]
\paragraph{Functors.} Given a $\mathcal{V}$-functor $G : \enriched{C} \to \enriched{D}$, $\tilde{F}(G) : \tilde{F}(\enriched{C}) \to \tilde{F}(\enriched{D})$ is given by the same mapping on objects and the following morphism on hom-objects
\[
(\tilde{F}\enriched{C})(A,B) = F(\enriched{C}(A,B)) \xrightarrow{F(G_{A,B})} F(\enriched{D}(GA,GB)) = (\tilde{F}\enriched{D})(GA,GB)
\]
\paragraph{Natural transformations.} Given a $\mathcal{V}$-natural transformation $\alpha : G \to H$ between $\mathcal{V}$-functors $G,H : \enriched{C} \to \enriched{D}$, $\tilde{F}(\alpha) : \tilde{F}(G) \to \tilde{F}(H)$ is given by the following family of morphisms
\[
\I_{\mathcal{W}} \xrightarrow{\cong} F(\I_{\mathcal{V}}) \xrightarrow{F(\alpha_{A})} F(\enriched{D}(GA,HA)) = (\tilde{F}\enriched{D})(GA,HA)
\]
\end{definition}

\begin{example}
    Let $\mathcal{C}$ be a usual $\catname{Set}$-category and let $F : \catname{Set} \to \catname{SLat}$ be the free semilattice functor.
    Then $\enriched{C} = \tilde{F}(\mathcal{C})$ has the same objects as $\mathcal{C}$ and every hom-object $\enriched{C}(A,B)$ additionally contains $f \join g$ for all $f,g \in \mathcal{C}(A,B)$ quotiented by the axioms of a semilattice, or, equivalently, consists of all finite non-empty subsets of $\mathcal{C}(A,B)$.
    The identity $\id_{A}$ is given by the identity of $\mathcal{C}(A,A)$ and the composition is defined componentwise.
    $\tilde{F}$ being a 2-functor means that $\enriched{C}$ inherits all adjunctions from $\mathcal{C}$.

    When the functor $\tilde{F}$ is induced by a free functor from $\catname{Set}$, we call it a free enrichment.
\end{example}

\begin{proposition}[Theorem 5.7.1~\cite{cruttwell2008normed}]
    Let $\mathcal{V}$ and $\mathcal{W}$ be symmetric monoidal categories and let $F : \mathcal{V} \to \mathcal{W}$ be a symmetric monoidal functor.
    Then, the induced 2-functor $\tilde{F} : \mathcal{V}\text{-}\catname{Cat} \to \mathcal{W}\text{-}\catname{Cat}$ maps monoidal $\mathcal{V}$-categories to monoidal $\mathcal{W}$-categories.
\end{proposition}

All the above properties are essential to give a semantics to (generalised) e-graphs.
The latter are defined over a signature of generating morphisms that induce some free category with, for example, symmetric monoidal closed structure.
Then, free enrichment allows us to talk about equivalence classes of morphisms in this category with all structures of the original category preserved.
We recap this interpretation next.

\subsection{Proofs of background results}
\label{app:background-proofs}

\begin{proof}[Proof of \Cref{prop:slat_closed_smc}]
    The tensor product of two semilattices $S_{1}$ and $S_{2}$ is defined as follows.
    $S_{1} \otimes S_{2}$ consists of pairs $(s_1,s_2)$ $s_{1} \in S_{1}$, $s_{2} \in S_{2}$ quotiented by commutativity, idempotence and associativity and additionally by the following relations
    \begin{itemize}
      \item $(s_{1} \join_{S_{1}} s_{1}',s_{2}) \equiv (s_{1},s_{2}) \join_{S_{1} \otimes S_{2}} (s_{1}',s_{2})$
      \item $(s_{1}, s_{2} \join_{S_{2}} s_{2}') \equiv (s_{1},s_{2}) \join_{S_{1} \otimes S_{2}} (s_{1},s_{2}')$
    \end{itemize}
  
    The unit for this tensor product is $I = \{*\}$---a one-element semilattice.
    Clearly $S \otimes I \cong S$ by mapping $(s,*) \mapsto s$ for all $s \in S$ and vice versa.
    The symmetry, associators and unitors are then obvious morphisms.
    Finally, the category is closed since the set of homomorphisms between two semilattices is a semilattice by defining $(f \join g)(x)$ as $f(x) \join g(x)$.
    The map $f \join g$ is a homomorphism since
    \[
    (f \join g)(x\join y)
    =f(x)\join f(y)\join g(x)\join g(y)
    =f(x)\join g(x)\join f(y)\join g(y)
    =(f\join g)(x)\join(f\join g)(y).
    \]
  \end{proof}

\begin{proof}[Proof of \Cref{prop:free_slat}]
	The functor $F$ is defined by letting $F(A) = \coprod_{A} I_{\catname{SLat}}$ where the latter is a coproduct which is a `free' product of semilattices defined as follows.
	The elements of $S_{1} \coprod S_{2}$ are expressions $s_{1}\join s_{2}\join\cdots\join s_{n}$ where each $s_{i}$ is either from $S_{1}$ or $S_{2}$, quotiented by all relations in $S_{1}$ and $S_{2}$.
	The adjunction is then given by the following natural isomorphisms
	\begin{align*}
		\catname{SLat}(\coprod_{A}(I_{\catname{SLat}}), B) & \cong \prod_{A}(\catname{SLat}(I_{\catname{SLat}}, B))       \\
		                                                   & \cong \catname{Set}(A,\catname{SLat}(I_{\catname{SLat}}, B)) \\
		                                                   & \cong \catname{Set}(A, U(B))
	\end{align*}
	The fist isomorphism is given by the fact that hom-functor makes limits into colimits in its first argument, the second being given by the property that $|\catname{Set}(A,B)| = |B|^{|A|} = |\underbrace{B \times \ldots \times B}_{|A|}|$, and the last isomorphism is given by the definition of $U$.
	Furthermore, we have,
	\[
		F(I) \cong I_{\catname{SLat}}
	\]
	and
	\begin{align*}
		F(A) \otimes F(B) & \cong (\coprod_{A} I_{\catname{SLat}}) \otimes (\coprod_{B} I_{\catname{SLat}}) \\
		                  & \cong \coprod_{A} (I_{\catname{SLat}} \otimes \coprod_{B} I_{\catname{SLat}})   \\
		                  & \cong \coprod_{A} (\coprod_{B} (I_{\catname{SLat}} \otimes I_{\catname{SLat}})) \\
		                  & \cong \coprod_{A} (\coprod_{B} I_{\catname{SLat}})                              \\
		                  & \cong \coprod_{A \times B} I_{\catname{SLat}}                                   \\
		                  & \cong F(A \times B)
	\end{align*}
	In the above we used the fact that functors $- \otimes X$ and $X \otimes -$ preserve colimits as they are both left-adjoint by symmetry and monoidal closedness of $\catname{SLat}$.
\end{proof}
